\section{Dual-Driver Relational Model and Persistence}

\subsection{Relational Schema Design}
Adhering to Codd's foundational relational principles (\cite{codd1970relational}), the EIMS data model normalizes entities into four core tables:
\begin{enumerate}
    \item \texttt{tenants}: Encapsulates enterprise organization entities, unique subdomains, and operational statuses.
    \item \texttt{roles} \& \texttt{permissions}: Mathematical bitmask definitions governing system privileges.
    \item \texttt{users}: Scoped user credentials, argon2id/bcrypt password hashes, and tenant references.
    \item \texttt{audit\_logs}: Append-only security audit trail recording actor IDs, actions, entity references, and IP addresses.
\end{enumerate}

\subsection{Automated Failover Relational Architecture}
In defense against container sabotage or network partition between application and database containers, \texttt{App\textbackslash Core\textbackslash Database} wraps PDO instantiation with dynamic failover:

\begin{lstlisting}[language=PHP, caption={Dual-Driver Connection Failover}]
try {
    $dsn = "mysql:host={$host};port={$port};dbname={$dbname};charset=utf8mb4";
    self::$pdo = new PDO($dsn, $user, $pass, [
        PDO::ATTR_ERRMODE => PDO::ERRMODE_EXCEPTION,
        PDO::ATTR_DEFAULT_FETCH_MODE => PDO::FETCH_ASSOC,
        PDO::ATTR_EMULATE_PREPARES => false,
    ]);
    return self::$pdo;
} catch (PDOException $e) {
    error_log('[WARN] Primary MySQL failed. Falling back to SQLite3.');
    return self::initializeSqliteFallback();
}
\end{lstlisting}

This abstraction guarantees that standard queries executed via \texttt{Database::executeQuery()} remain functional regardless of database backend state.
